% §10.2 CountMin
%
% Cross-reference resolution status:
%
% DEFINED here:
%   \label{sec:countmin}                          -- this section
%   \label{def:countmin}                          -- the construction
%   \label{alg:countmin-update-query}             -- update/query pseudocode
%   \label{thm:countmin-bound}                    -- the headline error bound
%
% RESOLVED (chapter 5):
%   \Cref{ch:approximate-sets}                    -- chapter 5
%   \Cref{def:bloom-filter}                       -- §5.2
%
% RESOLVED (chapter 9):
%   \Cref{ch:approximate-functions}               -- chapter 9
%   \Cref{def:n-ary-channel}                      -- §9.2
%   \Cref{thm:n-ary-kronecker}                    -- §9.3
%
% RESOLVED (chapter 10, prior section):
%   \Cref{sec:sketches-as-bernoulli}              -- §10.1
%
% FORWARD REFS within chapter 10 (resolve in Plan 10 Tasks 4-7):
%   \Cref{sec:hyperloglog}                        -- §10.3
%
% BIBLIOGRAPHIC (resolved in §10.6, Plan 10 Task 7):
%   cormode-muthukrishnan-2005                    -- CountMin original

\section{CountMin}
\label{sec:countmin}

\subsection*{Problem setup}

A data stream presents items $x_1, x_2, \ldots, x_n$ from a universe
$\Universe$ in some order, with possible repetition. For an item $a
\in \Universe$, write $f_a = |\{i : x_i = a\}|$ for its true frequency
in the stream, and write $F = \sum_b f_b = n$ for the total stream
length. The frequency-estimation query asks, for a given $a$, to
report an estimate $\hat f_a$ of $f_a$. Storing the exact frequencies
costs space proportional to the number of distinct items, which in
the heavy-hitter regime that motivates streaming sketches can run to
the billions. CountMin gives up exactness and returns a small
data structure of size polylogarithmic in the stream length, with a
query that runs in time independent of $F$ and an error bound on
$\hat f_a - f_a$ controllable by two design parameters.

\subsection*{The construction}

Fix a width $w$ and a depth $d$, both positive integers, and pick $d$
hash functions
\[
  h_1, h_2, \ldots, h_d \;:\; \Universe \longrightarrow \{1, 2, \ldots, w\}.
\]
The $h_i$ are drawn from a pairwise-independent family: for any two
distinct items $a, b \in \Universe$ and any two cells $u, v \in \{1,
\ldots, w\}$, $\Prob{h_i(a) = u \text{ and } h_i(b) = v} = 1/w^2$.
Pairwise independence is enough for the analysis below, and it is
much cheaper to construct than the full independence assumed for the
Bloom filter in \Cref{def:bloom-filter}.

The data structure is a $d \times w$ integer table $C$ with rows
indexed by $i \in \{1, \ldots, d\}$ and columns indexed by $j \in
\{1, \ldots, w\}$, initialized to all zeros.

\begin{definition}[CountMin sketch]
\label{def:countmin}
The \emph{CountMin sketch} of parameters $(w, d)$ is the data
structure $(C, h_1, \ldots, h_d)$ together with the two operations:
\begin{itemize}
  \item \textsc{Update}$(a)$: for $i = 1, \ldots, d$, set
    $C[i][h_i(a)] \gets C[i][h_i(a)] + 1$.
  \item \textsc{Query}$(a)$: return $\hat f_a = \min_i C[i][h_i(a)]$.
\end{itemize}
After processing a stream by calling \textsc{Update} once per item,
the query returns the estimated frequency of any item $a$.
\end{definition}

\Cref{alg:countmin-update-query} states the two operations in
pseudocode.

\begin{algorithm}[ht]
  \caption{CountMin update and query}
  \label{alg:countmin-update-query}
  \SetKwInOut{Input}{input}
  \SetKwInOut{Output}{output}
  \SetKwProg{func}{procedure}{}{}
  \SetKwProg{fn}{function}{}{}

  \func{\textsc{Update}($a$)}{
    \For{$i \gets 1$ \KwTo $d$}{
      $C[i][h_i(a)] \gets C[i][h_i(a)] + 1$\;
    }
  }
  \BlankLine
  \fn{\textsc{Query}($a$)}{
    $m \gets \infty$\;
    \For{$i \gets 1$ \KwTo $d$}{
      \lIf{$C[i][h_i(a)] < m$}{$m \gets C[i][h_i(a)]$}
    }
    \Return $m$\;
  }
\end{algorithm}

Both operations run in $O(d)$ time per call, independent of the
stream length $F$ and the universe size $|\Universe|$. The data
structure uses $d \cdot w$ counters of storage, also independent of
$F$. The minimum across rows in \textsc{Query} is the central trick:
any single row $i$ records a count $C[i][h_i(a)]$ that is at least
$f_a$ (one increment per occurrence of $a$ lands in that cell), plus
the contributions of any other items whose hash $h_i$ collides with
$a$'s. Taking the minimum across $d$ independent rows attenuates the
collision noise: a row only inflates the estimate if some heavy
collider hashes to the same cell, and the rows are independent enough
that this happens in all $d$ rows only with small probability.

\subsection*{Why the estimate never under-shoots}

For every row $i$, the cell $C[i][h_i(a)]$ accumulates a count of
$f_a$ from $a$'s own updates, plus a non-negative count from any
other item $b \ne a$ whose hash $h_i(b)$ equals $h_i(a)$. Both
contributions are non-negative, so $C[i][h_i(a)] \ge f_a$ for every
$i$, and therefore
\[
  \hat f_a \;=\; \min_i C[i][h_i(a)] \;\ge\; f_a.
\]
The estimate is always at least the true count. This is the
deterministic counterpart of the Bloom filter's deterministic
no-false-negatives property (\Cref{def:bloom-filter}): the only
failure mode is over-estimation, and the error is one-sided. The
analysis below bounds the size of the over-estimate in probability.

\subsection*{The error bound}

\begin{theorem}[CountMin bound]
\label{thm:countmin-bound}
Fix $\epsilon, \delta \in (0, 1)$, and run CountMin with width
$w = \lceil e/\epsilon \rceil$ and depth $d = \lceil \ln(1/\delta) \rceil$.
For any item $a \in \Universe$, the estimate $\hat f_a$ returned by
\textsc{Query}$(a)$ satisfies
\[
  f_a \;\le\; \hat f_a \;\le\; f_a + \epsilon F
\]
with probability at least $1 - \delta$, where $F = \sum_b f_b$ is
the total stream length.
\end{theorem}

\begin{proof}
Fix the query item $a$ and a single row $i \in \{1, \ldots, d\}$.
Decompose the cell $C[i][h_i(a)]$ at the end of the stream into the
contribution from $a$ itself and the contribution from collisions:
\[
  C[i][h_i(a)]
  \;=\; f_a \;+\; \sum_{b \ne a} f_b \cdot \mathbf{1}[h_i(b) = h_i(a)],
\]
where $\mathbf{1}[\cdot]$ is the indicator. Write $X_i = \sum_{b \ne
a} f_b \cdot \mathbf{1}[h_i(b) = h_i(a)]$ for the row-$i$ collision
mass, so that $C[i][h_i(a)] = f_a + X_i$ and $X_i \ge 0$.

By pairwise independence of $h_i$, the marginal probability that any
specific $b \ne a$ collides with $a$ under $h_i$ is $\Prob{h_i(b) =
h_i(a)} = 1/w$. Linearity of expectation gives
\[
  \mathbb{E}[X_i]
  \;=\; \sum_{b \ne a} f_b \cdot \Prob{h_i(b) = h_i(a)}
  \;=\; \frac{F - f_a}{w}
  \;\le\; \frac{F}{w}.
\]
Apply Markov's inequality to the non-negative random variable $X_i$
with threshold $\epsilon F$:
\[
  \Prob{X_i \ge \epsilon F}
  \;\le\; \frac{\mathbb{E}[X_i]}{\epsilon F}
  \;\le\; \frac{1}{\epsilon w}.
\]
With $w = \lceil e/\epsilon \rceil \ge e/\epsilon$, this bound becomes
\[
  \Prob{X_i \ge \epsilon F} \;\le\; \frac{1}{e}.
\]
The events $\{X_i \ge \epsilon F\}$ for $i = 1, \ldots, d$ depend on
independently drawn hash functions $h_1, \ldots, h_d$, so they are
mutually independent. The query reports the minimum across rows, and
$\hat f_a - f_a > \epsilon F$ exactly when \emph{every} row
over-shoots:
\[
  \Prob{\hat f_a > f_a + \epsilon F}
  \;=\; \Prob{X_i \ge \epsilon F \text{ for all } i}
  \;\le\; (1/e)^d.
\]
With $d = \lceil \ln(1/\delta) \rceil \ge \ln(1/\delta)$,
$(1/e)^d \le \delta$, which is the claimed bound. The lower bound
$\hat f_a \ge f_a$ is the deterministic argument of the previous
subsection. \qedhere
\end{proof}

Two design points are worth flagging. First, the bound is additive
in $F$, not in $f_a$: heavy hitters with $f_a$ comparable to $F$ are
estimated with small \emph{relative} error, while light items with
$f_a \ll \epsilon F$ may be reported as $\epsilon F$ and the relative
error is essentially uninformative. CountMin is the right sketch when
the question of interest is which items are heavy, not what the
counts of the light items are. Second, the space cost is
\[
  d \cdot w \;=\; \lceil \ln(1/\delta) \rceil \cdot \lceil e/\epsilon \rceil
  \;=\; O\bigl((1/\epsilon) \ln(1/\delta)\bigr)
\]
counters, which is logarithmic in the failure probability and inverse
in the additive error parameter. The original analysis is due to
\textcite{cormode-muthukrishnan-2005}; the proof above is the
standard Markov-plus-union-bound presentation.

\subsection*{The $\Bernoulli{T}$ type}

CountMin instantiates the framework of \Cref{ch:approximate-functions}
with output type $T = \mathbb{N}$. The query is a function from
streams to non-negative integers, and its error behaviour is
summarized by a channel matrix $Q$ in the sense of
\Cref{def:n-ary-channel}: rows are indexed by the true count $f_a$
and columns by the reported estimate $\hat f_a$, with entry
\[
  Q_{f_a, \hat f_a} \;=\;
  \Prob{\textsc{Query}(a) = \hat f_a \Given \text{true count is } f_a}.
\]
The deterministic lower bound $\hat f_a \ge f_a$ has a clean
expression in $Q$: every row places all its mass on columns at or
above the diagonal, so $Q$ is upper triangular. This is the analogue,
for a count-valued output, of the upper-triangular channel matrix
associated with the Bloom filter as a one-sided approximate set: the
deterministic lower bound on the estimate corresponds structurally to
the deterministic ``no false negatives'' guarantee of the membership
report. The probabilistic upper bound from
\Cref{thm:countmin-bound} fills in the row-wise shape: each row's
mass is concentrated near the diagonal, with at most a $\delta$
fraction of probability sitting more than $\epsilon F$ columns to the
right of $f_a$. The matrix is large (rows and columns indexed by all
of $\mathbb{N}$, in principle), but its structure is summarized by
the two parameters $(\epsilon, \delta)$.

\subsection*{Composition: merging two CountMin sketches}

A common operational pattern is that two CountMin sketches $C_1$ and
$C_2$ are built independently on two disjoint substreams of a single
larger stream, and the analysis wants a sketch of the combined
stream. If both sketches use the same parameters $(w, d)$ and the
same hash functions $h_1, \ldots, h_d$, the merge is entry-wise
addition:
\[
  (C_1 \oplus C_2)[i][j] \;=\; C_1[i][j] + C_2[i][j]
  \quad \text{for all } i \in \{1, \ldots, d\}, \; j \in \{1, \ldots, w\}.
\]
The merged table is exactly the CountMin sketch that would have been
built by running \textsc{Update} on the concatenation of the two
substreams, because \textsc{Update} only increments counters and
addition is commutative. \textsc{Query} on the merged sketch returns
$\min_i (C_1 + C_2)[i][h_i(a)]$, which is the estimate of the total
count $f_a^{(1)} + f_a^{(2)}$ in the combined stream. The same
$(\epsilon, \delta)$ bound from \Cref{thm:countmin-bound} applies,
with the new $F$ being the sum of the two substream lengths.

In the language of \Cref{ch:approximate-functions}, this merge is a
$\Bernoulli{\mathbb{N}}$ composition. The channel matrix of the
merged sketch is the channel matrix that
\Cref{thm:n-ary-kronecker} would assign to the joint query of the
two substream sketches under the natural product of integer counts.
Crucially the merge does not double the failure probability: the two
sketches share their hashes, so the row-collision events on the
merged sketch are the same events as on each constituent, not the
union of two separate event sets. Sharing the hash family is the
discipline that lets independently built sketches add up to one
correct sketch of the union.

\subsection*{A worked example}

Take $w = 10$ and $d = 5$, a small but illustrative configuration.
The corresponding error parameters from \Cref{thm:countmin-bound} are
$\epsilon \le e/w \approx 0.272$ (the additive bound is at most
$0.272 \cdot F$) and $\delta \le e^{-d} \approx 0.00674$ (the
failure probability is at most $0.67\%$).

Consider a stream of $F = 1000$ items drawn from a universe of $50$
distinct items, with frequency profile
\[
  f_{a_1} = 400, \; f_{a_2} = 200, \; f_{a_3} = 100, \;
  f_{a_4} = \cdots = f_{a_{20}} = 10, \;
  f_{a_{21}} = \cdots = f_{a_{50}} = 4,
\]
which sums to $400 + 200 + 100 + 17 \cdot 10 + 30 \cdot 4 = 990$,
and pad with $10$ further occurrences of $a_1$ to reach exactly
$F = 1000$ (so $f_{a_1} = 410$). Query for the heavy hitter $a_1$
with $f_{a_1} = 410$. The expected per-row over-estimate is
\[
  \mathbb{E}[X_i] \;=\; (F - f_{a_1})/w \;=\; 590/10 \;=\; 59,
\]
so each row reports a value near $410 + 59 = 469$ in expectation,
and the minimum across $d = 5$ rows is pulled below this mean by the
order statistic. A typical realization (verified in the chapter
notebook with seed 0) returns $\hat f_{a_1} = 432$, an over-estimate
of $22$, well inside the headline bound $\epsilon F = 0.272 \cdot
1000 = 272$. Now query for a light item $a_{30}$ with $f_{a_{30}} =
4$. The expected per-row over-estimate is $(1000 - 4)/10 = 99.6$, so
even the minimum across five rows will likely report something on
the order of $50$ to $90$. The estimate is correctly bounded above by
$f_{a_{30}} + \epsilon F = 4 + 272 = 276$, but it is essentially
uninformative about the true count of $4$: the additive error term
is twenty times the true frequency. The headline bound holds, but it
holds in a regime where light items are washed out by the collision
mass of heavier ones. This is exactly the heavy-hitter design point
of CountMin: heavy items are estimated tightly relative to their own
count, and light items are reported only up to the additive $\epsilon
F$ slack.

\subsection*{Looking ahead}

CountMin estimates the frequency of a specific item. The next
section, \Cref{sec:hyperloglog}, turns to a complementary streaming
question: how many \emph{distinct} items have been seen, regardless
of which ones. The answer is not a per-item count but a single
integer, and the sketch that estimates it works on a completely
different principle, the leading-zero count of a random hash, with a
space cost that is logarithmic in the cardinality being estimated
rather than logarithmic in the failure probability.
